% Generated from locale/tr/content/counterfactuals/minimal-change-semantics/introduction.tex; do not hand-edit.


\olfileid[tr]{cnt}{min}{int}

\olsection{Giriş}

Stalnaker ile Lewis ``Kibrit çakılsaydı yanardı'' gibi karşıolgusal koşullar
için açıklamalar önerdi. Bu açıklamalar, böyle tümcelerin doğruluk
koşullarının nasıl doğru anlaşılması gerektiğine ilişkin önerilerdi. İki
önerinin de ardındaki düşünce şudur: Bir karşıolgusal koşulun doğru olup
olmadığını değerlendirmek için, önbileşeni doğru kılmak üzere gerçek dünyanın
durumundan en az ölçüde ayrılan olası dünyaları incelemeliyiz. Bu olası
dünyalarda artbileşen doğruysa karşıolgusal doğrudur. Örneğin elimde bir
kibritle kibrit kutusu tuttuğumu varsayalım. Gerçek dünyada yalnızca onlara
bakıyor ve kibriti çaksaydım ne olacağını düşünüyorum. Kibriti çaktığım, fakat
gerçek dünyadan en az ayrılan durumda eyleme geçmeye karar verir ve kibriti
çakarım. Başka hiçbir şeyin değişmemesi anlamında bu ayrım en azdır: Aynı
zamanda havaya sıçramam, kibriti çakmam saçımı da tutuşturmaz, parmaklarımdaki
bütün gücü birden yitirmem, su tabancasıyla kurulmuş bir pusuda aynı anda
sırılsıklam edilmem vb. Bu alternatif olasılıkta kibrit yanar. Dolayısıyla
kibriti çaksaydım yanacağı doğrudur.

Bu sezgisel açıklama, karşıolgusal mantıkların biçimsel semantiğiyle
birleştirilebilir. Lewis karşıolgusal için ``$\boxright$'' simgesini,
Stalnaker ise~``$>$'' simgesini kullandı. Biz~$\cif$'i kullanacak ve onu
önerme mantığına ikili bağlaç olarak ekleyeceğiz. Böylece $!A\lif !B$
biçimindeki !!{formula}s{}in yanı sıra $!A\cif !B$ biçiminde !!{formula}s da
vardır. Modal mantığın bağıntısal semantiği gibi biçimsel semantik de
!!{formula}s{}in dünyalarda değerlendirildiği modellere dayanır;
$\mSat{M}{!A\cif !B}[w]$'yi tanımlayan sağlama koşulu, bazı (başka)
dünyalar~$w'$ için $\mSat{M}{!A}[w']$ ile $\mSat{M}{!B}[w']$ kullanılarak
verilir. Hangi $w'$? Sezgisel olarak $\mSat{M}{!A}[w']$ olan ve~$w$'ye en
yakın dünya ya da dünyalar. Bu nedenle modele bir ``yakınlık'' bağıntısının
da eklenmesi gerekir.

Lewis karşıolgusal durumları grafiksel olarak göstermenin öğretici bir
yolunu sundu. Her olası dünya, başka dünyaları içeren iç içe küreler
kümesinin merkezindedir---bu küreleri eş merkezli çemberler olarak çizeriz.
İki küre arasındaki dünyalar merkezdeki dünyaya birbirleri kadar yakındır;
içteki bir kürede bulunanlar daha yakın, dıştaki bir kürede bulunanlar daha
uzaktır.
\begin{center}
\begin{tikzpicture}[scale=.7]\small
  \spheresystem{5}
  \draw (0,0) node {$w$};
  \propositionintersect{0}{4}{40}{3.5}
  {\tikzset{proposition/.append={smooth,tension=1.4}}
  \proposition[shift={(1.6,-1)}]{90}{1}{30}{4}}
  \path (1.5,3) node[below] {$\formula{B}$};
  \path (3,0) node {$\formula{A}$};
\end{tikzpicture}
\end{center}
En yakın $!A$-dünyaları, $!A$'nın sağlandığı ve merkezdeki~$w$ dünyasının
çevresindeki en küçük kürede (gri alanda) bulunan $w'$ dünyalarıdır.
Sezgisel olarak $!A\cif !B$, en yakın bütün $!A$-dünyalarında $!B$ doğruysa
$w$'de sağlanır.

